\section{KL 数值：从确定性代理到无偏残差校正}
\label{sec:kl-scalar}
本节固定目标 $K=I_p(r)$，所以总定理中的 $f_v=r_v$ 是标量。随机估计量默认使用 $\rho=p$；任意 $\rho$ 的权重与支持条件见第~\ref{sec:kl-sampling} 节。先给出 Top-$J$ 记号，再分类确定性代理、全词表单样本控制和稀疏头尾控制。每个估计量的误差都相对同一个完整 $K$。
\subsection{头部与尾部记号}

沿用前节的 $S,P,Q,\tau_p,\tau_q,J_A,H$ 和 $\ell_{\tail}$。令
\[
 K=\KL(p\Vert q)=\sum_vp(v)r_v,
 \qquad H=\sum_{v\in S}p(v)r_v,
 \]
并令尾部条件分布为
\[
 p_\tail(v)=\frac{p(v)}{\tau_p},\qquad
 q_\tail(v)=\frac{q(v)}{\tau_q},\qquad v\notin S.
\]
定义尾部质量项和尾部形状项
\[
 M=\tau_p\log\frac{\tau_p}{\tau_q},\qquad
 B=\tau_p\KL(p_\tail\Vert q_\tail)\geq0.
\]

\begin{corollary}[头尾分区的 KL 链式分解]
\label{thm:chapter3-kl-chain}
在 $p\ll q$ 且 $0<\tau_p,\tau_q<1$ 时，
\begin{equation}
 K=H+M+B.
\label{eq:chapter3-kl-chain}
\end{equation}
\end{corollary}

\begin{proof}
在定理~\ref{thm:kl-partition-chain} 中取头部 singleton 与单个尾桶。singleton 的条件 KL 为零，粗粒度 KL 是 $H+M$，唯一剩余的条件形状项是 $B$，故得到所述恒等式。
\end{proof}

该分解说明 sparse teacher payload 能够知道 $H$、$P$、$Q$，因而知道 $M$，但不知道尾部形状项 $B$。所有 Top-$J$ scalar estimator 的差异都可以直接归结为是否丢弃、聚合或抽样估计 $B$。

\subsection{确定性估计量}

截断 Top-$J$ 估计量定义为
\[
 \widehat K^{\mathrm{trunc}}=H.
\]
由式~\eqref{eq:chapter3-kl-chain}，其条件偏差和方差为
\begin{equation}
 \operatorname{Bias}(\widehat K^{\mathrm{trunc}})=-M-B,
 \qquad
 \Var(\widehat K^{\mathrm{trunc}})=0.
 \label{eq:chapter3-trunc-bias-var}
\end{equation}
它删除了尾部贡献；偏差 $-M-B$ 可正可负，因为 $M$ 的符号可以正也可以负，而 $B$ 总是非负。

Tail-aware scalar 估计量把尾部聚合为一个伪 token：
\[
 \widehat K^{\mathrm{TA}}=H+M.
\]

\begin{proposition}[Tail-aware scalar 的偏差]
\label{prop:chapter3-ta-scalar}
\[
 \operatorname{Bias}(\widehat K^{\mathrm{TA}})=-B,
 \qquad
 \Var(\widehat K^{\mathrm{TA}})=0.
\]
因此它是完整 KL 的确定性下界，且当且仅当尾部条件分布相等时达到完整 KL。
\end{proposition}

\begin{proof}
将 $\widehat K^{\mathrm{TA}}=H+M$ 与 $K=H+M+B$ 相减即可。$B\geq0$ 来自 KL 非负性。
\end{proof}

Normalized Top-$J$ 不是完整 KL 的估计量，而是头部条件分布的另一个目标：
\[
 K^{\mathrm{norm}}_J
 =\KL\left(\frac{p|_S}{P}\middle\|\frac{q|_S}{Q}\right)
 =\frac{H}{P}+\log\frac{Q}{P}.
\]
它的条件方差为零，但相对完整 KL 的确定性偏差为
\begin{equation}
 \operatorname{Bias}(K^{\mathrm{norm}}_J)
 =\frac{H}{P}+\log\frac{Q}{P}-(H+M+B).
 \label{eq:chapter3-normalized-bias}
\end{equation}
这个偏差没有固定符号，因为 Normalized Top-$J$ 同时改变了头部权重并删除了尾部形状。

\subsection{单 token KL 估计器的基准族}

先记录不依赖 Top-$J$ 的单 token 族。令 $A\sim p$，$r=r_A$，并假设相应矩存在。反向 KL 的直接 estimator 为
\[
 \widehat K_1=r.
\]
它满足
\[
 \E[\widehat K_1]=K,\qquad
 \operatorname{Bias}(\widehat K_1)=0,\qquad
 \Var(\widehat K_1)=\Var_p(r).
\]
常用的平方 estimator 为
\[
 \widehat K_2=\frac12r^2,
\]
其相对于 $K$ 的偏差和方差为
\begin{equation}
 \operatorname{Bias}(\widehat K_2)=\frac12\E_p[r^2]-K,
 \qquad
 \Var(\widehat K_2)=\frac14\left(\E_p[r^4]-\E_p[r^2]^2\right).
 \label{eq:chapter3-k2-bias-var}
\end{equation}
非负修正 estimator 为
\[
 \widehat K_3=e^{-r}+r-1.
\]
因为 $e^{-r}=q(A)/p(A)$，有
\[
 \E_p[e^{-r}]=\sum_vq(v)=1,
\]
故 $\E[\widehat K_3]=K$，且
\begin{equation}
 \Var(\widehat K_3)=\Var_p(e^{-r}+r-1).
 \label{eq:chapter3-k3-var}
\end{equation}
K3 的方差还可展开为
\[
 \Var(\widehat K_3)
 =\Var_p(r)+\Var_p(e^{-r})+2\operatorname{Cov}_p(r,e^{-r}).
\]
因此 K3 比 K1 方差小的充要条件是
\[
 \Var_p(e^{-r})+2\operatorname{Cov}_p(r,e^{-r})<0.
\]
K3 的逐样本非负性来自 $e^{-r}\geq1-r$，并不推出全局方差排序。对 $q(v)=p(v)(1+\varepsilon u_v)$、$\E_p[u]=0$ 的局部扰动，逐项 Taylor 展开给出
\[
 r_v=-\varepsilon u_v+\frac{\varepsilon^2u_v^2}{2}+O(\varepsilon^3),
 \quad K=\frac{\varepsilon^2}{2}\E_p[u^2]+O(\varepsilon^3),
\]
\[
 \widehat K_2=\frac{\varepsilon^2u_A^2}{2}+O(\varepsilon^3),
 \quad \widehat K_3=\frac{\varepsilon^2u_A^2}{2}+O(\varepsilon^3).
\]
这解释了 K2 的局部二阶正确性，以及 K3 在两分布接近时将单点波动从一阶降到二阶的机制。
K1 是总定理中 $a=0$ 的特例。K3 则取 $a_v=1-e^{-r_v}$，因为 $I_p(a)=1-\sum_vq(v)=0$，总定理给出 $r_A-a_A=e^{-r_A}+r_A-1$。K2 改变了数值目标而未补回残差，所以需要直接报告上面的偏差。

\subsection{采样尾部的无偏 scalar 估计量}

令
\[
 \mu_r=\E[r_A\mid A\notin S]
 =\log\frac{\tau_p}{\tau_q}+\KL(p_\tail\Vert q_\tail).
\]
则 $\tau_p\mu_r=M+B$。头部精确、尾部采样（Head-Exact Tail-Sampled，HETS）的数值估计量为
\begin{equation}
 \widehat K^{\mathrm{HETS}}=H+J_A r_A.
 \label{eq:chapter3-ema-scalar}
\end{equation}
它是总定理中头部 $a_v=r_v$、尾部 $a_v=0$ 的特例。本章后续统一称为 HETS 数值估计量；EMA 只在对应原项目名称时出现。时间上的指数滑动平均不属于此命名。

\begin{theorem}[HETS 数值的无偏性与方差]
\label{thm:chapter3-ema-scalar}
若 $A\sim p$ 且 $\E_p[r_A^2]<\infty$，则
\[
 \E[\widehat K^{\mathrm{HETS}}]=K,
\]
并且
\begin{align}
 \Var(\widehat K^{\mathrm{HETS}})
 &=\tau_p\Var(r_A\mid A\notin S)+P\tau_p\mu_r^2\\
 &=\sum_{v\notin S}p(v)r_v^2
   -\left(\sum_{v\notin S}p(v)r_v\right)^2.
 \label{eq:chapter3-ema-scalar-var}
\end{align}
\end{theorem}

\begin{proof}
尾部修正的期望为
\[
 \E[J_A r_A]=\sum_{v\notin S}p(v)r_v=\tau_p\mu_r=M+B.
\]
加上 $H$ 得到 $K$。对随机变量 $J_A r_A$ 按 $J_A$ 使用全方差公式：头部条件下它为零，尾部条件下它为 $r_A$，因此
\[
 \Var(J_A r_A)=\tau_p\Var(r_A\mid A\notin S)+P\tau_p\mu_r^2.
\]
\end{proof}

还可以用尾部桶 scalar
\[
 \ell_{\tail}=\log\frac{\tau_p}{\tau_q}
\]
作残差化：
\begin{equation}
 \widehat K^{\mathrm{Res\text{-}scalar}}
 =H+\tau_p\ell_{\tail}+J_A(r_A-\ell_{\tail}).
 \label{eq:chapter3-res-scalar}
\end{equation}
它满足
\[
 \E[\widehat K^{\mathrm{Res\text{-}scalar}}]=K,
\]
且
\begin{equation}
 \Var(\widehat K^{\mathrm{Res\text{-}scalar}})
 =\tau_p\Var(r_A\mid A\notin S)+P\tau_p(\mu_r-\ell_{\tail})^2.
 \label{eq:chapter3-res-scalar-var}
\end{equation}
所以 scalar 残差化相对 HETS 的方差差为
\begin{equation}
 \Var(\widehat K^{\mathrm{HETS}})-\Var(\widehat K^{\mathrm{Res\text{-}scalar}})
 =P\tau_p(2\mu_r\ell_{\tail}-\ell_{\tail}^2).
 \label{eq:chapter3-res-scalar-gap}
\end{equation}
它严格降低方差当且仅当 $|\mu_r-\ell_{\tail}|<|\mu_r|$。当固定的 $\ell_{\tail}\ne0$ 且尾部内部形状 KL 足够小时，$\mu_r$ 足够接近 $\ell_{\tail}$，严格降方差条件成立；若 $\ell_{\tail}=0$，两估计量逐样本相同。

\begin{proposition}[标量尾部控制族]
对任意在抽样前固定的实数 $\beta$，令
\[
 \widehat K(\beta)=H+\tau_p\beta+J_A(r_A-\beta).
\]
则它对 $K$ 无偏，且
\[
 \Var(\widehat K(\beta))
 =\tau_p\Var_{\tail}(r_A)+P\tau_p(\mu_r-\beta)^2.
\]
因此最优标量控制为 $\beta^*=\mu_r$；HETS 取 $\beta=0$，式~\eqref{eq:chapter3-res-scalar} 取 $\beta=\ell_{\tail}$。
\end{proposition}

\begin{proof}
先计算均值：
\[
 \E[\widehat K(\beta)]=H+\tau_p\beta+\tau_p(\mu_r-\beta)=H+\tau_p\mu_r=K.
\]
再计算中心化随机项的二阶矩：
\begin{align*}
 \Var(\widehat K(\beta))
 &=\E[J_A(r_A-\beta)^2]-\{\E[J_A(r_A-\beta)]\}^2\\
 &=\tau_p\{\Var_{\tail}(r_A)+(\mu_r-\beta)^2\}
   -\tau_p^2(\mu_r-\beta)^2\\
 &=\tau_p\Var_{\tail}(r_A)+P\tau_p(\mu_r-\beta)^2.
\end{align*}
最后一项非负并在 $\beta=\mu_r$ 时为零，给出最优值。将 $\beta=0$ 与 $\beta=\ell_{\tail}$ 代入并相减即得式~\eqref{eq:chapter3-res-scalar-gap}。
\end{proof}

这个 scalar 构造是本章从控制族导出的 KL 数值估计量，应与论文中以向量梯度为目标的 ResOPD 算法区别命名。因为 $\mu_r=\ell_{\tail}+D_{\tail}$，其中 $D_{\tail}=\KL(p_\tail\Vert q_\tail)$，方差改变量也等于
\[
 P\tau_p\{\ell_{\tail}^2+2\ell_{\tail}D_{\tail}\}.
\]
当 $\ell_{\tail}\geq0$ 时它非负；当 $\ell_{\tail}<0$ 时，它在 $D_{\tail}>-\ell_{\tail}/2$ 下为负。这给出 scalar 残差化反而增大方差的明确条件。
